\section{敏感性分析}
\label{sec:sensitivity}

本节从三个维度检验主结论的稳健性：第一问中 small adenoma 样本的归类方案、第二问中稳定特征筛选阈值 $\tau$、第三问中跨疾病应用的阈值偏移。三处敏感性分析要么验证主结论稳健，要么如实报告模型的适用边界。

\subsection{第一问：small adenoma 四种方案敏感性（Zeller）}

Zeller CRC 数据集中存在一类 small adenoma（小腺瘤）样本，其疾病状态介于健康与癌症之间。主分析将其归入健康组（方案 (1)），但这一选择可能影响模型性能。测试了四种归类方案：(1)归健康（默认主方案）、(2)归病变（与 cancer 合并为正类）、(3)剔除（仅保留明确的健康和 cancer 样本）、(4)单开一类（二分类问题中部分等价于剔除）。

结果如下表（\cref{tab:adenoma}）。方案 (1)的 L2 AUC 为 0.7907，RF 为 0.8454（$n=121$）。方案 (2)将 small adenoma 归为病变后，性能骤降至 L2 0.6112、RF 0.6509——小腺瘤的微生物组特征更接近健康而非癌症，将其混入正类会严重污染标签，导致模型性能下降，因此方案 (2)被排除。方案 (3)剔除 small adenoma 后，L2 AUC 为 0.8022、RF 为 0.8667（$n=95$），相对方案 (1)的变化分别为 +0.0115 和 +0.0213，均小于 0.05 的显著性阈值，说明主方案 (1)的选择对结果影响不大。方案 (4)单开一类在二分类框架下与方案 (3)结果相同。

综合来看，主方案 (1)（归健康）是稳健的选择，方案 (3)和(4)的结果作为敏感性附录留存。

\textbf{small adenoma 四种方案敏感性表}（Zeller）：

{\small
\begin{longtable}{>{\centering\arraybackslash}p{0.3800\textwidth}>{\centering\arraybackslash}p{0.1800\textwidth}>{\centering\arraybackslash}p{0.1600\textwidth}>{\centering\arraybackslash}p{0.1600\textwidth}}
\caption{small adenoma 四种方案敏感性表（Zeller）}\label{tab:adenoma}\\
\toprule
方案 & L2 AUC & RF AUC & $n$ \\
\midrule
\endfirsthead
\toprule
方案 & L2 AUC & RF AUC & $n$ \\
\midrule
\endhead
\bottomrule
\endlastfoot
(1) 归健康 & 0.7907 & 0.8454 & 121 \\
(2) 归病变 & 0.6112 & 0.6509 & 121 \\
(3) 剔除 & 0.8022 & 0.8667 & 95 \\
(4) 单开一类 & 0.8022 & 0.8667 & 95 \\
\end{longtable}
}

\subsection{\texorpdfstring{第二问：$\tau$ 敏感性（全量 bootstrap）}{第二问：tau 敏感性（全量 bootstrap）}}

第二问中稳定生物标志物的筛选依赖于阈值 $\tau$——一个特征在 bootstrap 重采样中被选中的频率超过 $\tau$ 才被认定为「稳定」。$\tau$ 的取值直接影响入选特征的数量，测试了 $\tau=0.4$、$\tau=0.5$、$\tau=0.6$、$\tau=0.7$ 四个水平。

入选稳定特征数随 $\tau$ 单调下降：CRC 从 6 个降至 2 个，IBD 从 6 个降至 2 个，Obesity 从 32 个降至 3 个。$\tau=0.5$ 是稳定簇与噪声长尾的自然分界点——低于 0.5 时入选数快速膨胀（大量低频率特征涌入），高于 0.5 时入选数趋于平稳。主分析采用 $\tau=0.5$，在该阈值下 CRC 和 IBD 各有 4 个稳定标志物，Obesity 有 20 个。

无论 $\tau$ 如何变化，主结论始终成立：CRC 和 IBD 有少量高度稳定的标志物（$\tau=0.7$ 时仍有 2 个），而 Obesity 的标志物数量对 $\tau$ 高度敏感（从 32 骤降至 3），反映其信号弱、稳定性差的本质。$\tau$ 敏感性分析验证了稳定标志物选择的稳健性。

\subsection{第三问：阈值偏移（C3 组合）}

{\raggedright
跨疾病应用中，训练集确定的决策阈值迁移到测试集时可能因患病基线不同而失效。以 C3 组合（训练 CRC+IBD，测试 Obesity）为例量化阈值偏移的严重程度。\par}

C3 训练集的患病基线为 0.3160，测试集为 0.6482，基线差 $\Delta=+0.3322$——测试集的正类比例是训练集的两倍以上。在训练集上由 Youden J 确定的最优阈值 $\tau^{*}=0.9205$，落到测试集分数分布的 96.0\% 分位，即只有 4\% 的测试样本分数超过该阈值，灵敏度仅 0.0244，模型几乎将全部测试样本判为健康。

Platt 校准是单调变换，不改变样本的相对排序和决策边界，因此无法修复这种基线偏移导致的阈值偏移——校准前后 Youden 阈值迁移的灵敏度一致。即便放弃训练集阈值、改用 0.5 的朴素概率阈值，灵敏度也仅 0.1646，仍然很低。C3 的案例证伪了「只要 AUC 尚可就能实际应用」的假设——当测试集患病基线与训练集差异过大时，训练集拟合的任何单调校准都无法给出合理的决策阈值。阈值偏移是跨疾病应用中独立于 AUC 的另一重障碍，诊断图见模型构建节 \cref{fig:s3-threshold-drift}。
\fussy

\subsection{讨论}

三处敏感性分析从不同角度刻画了模型的可靠性与边界。

第一问的 small adenoma 四种方案分析表明，主分析将小腺瘤归为健康组的选择是稳健的——剔除后性能变化小于 0.05，而归为病变则会因标签污染导致性能骤降。这说明小腺瘤的微生物组特征更接近健康状态，在 CRC 预测中不应被视为阳性样本。

第二问的 $\tau$ 敏感性验证了稳定标志物筛选的可靠性。CRC 和 IBD 的标志物在高阈值下仍能保留，说明这些特征确实具有高度的重采样稳定性；而 Obesity 的标志物数量随 $\tau$ 急剧下降，进一步印证了其弱信号本质。$\tau=0.5$ 作为稳定簇与噪声长尾的分界点是合理的选择。

第三问的阈值偏移分析则揭示了跨疾病预测的一个关键边界：即使模型的排序能力（AUC）尚可，训练集与测试集之间的患病基线差异也会导致决策阈值完全失效。这一发现对实际应用有直接意义——跨疾病应用不能直接迁移训练集阈值，必须针对目标疾病的患病基线做无监督的阈值重定，或采用不依赖绝对阈值的排序输出。

总体而言，第一问和第二问的主结论在敏感性测试下保持稳健，第三问的敏感性分析则如实暴露了跨疾病应用的适用边界。这些结果共同支撑了论文的核心发现：单疾病预测模型可靠且稳健，生物标志物筛选具有统计依据，而跨疾病预测在当前数据条件下仍面临信号特异和阈值偏移的双重障碍。
